\section*{Resumen}
\addcontentsline{toc}{section}{Resumen}

Esta segunda etapa parte del clasificador obtenido en la etapa anterior y se
concentra en los problemas que quedaron abiertos: el bajo rendimiento de
potasio, el desbalance, el sesgo laboratorio/campo, la selección del modelo y
una evaluación final que no reutilizara el test histórico. Se reconstruyó la
publicación ampliada del dataset, se verificaron 33\,437 archivos y se
eliminaron cuatro duplicados exactos. El corpus definitivo contiene 33\,433
imágenes de nueve clases, con un holdout nuevo de 5\,015 fotografías congelado
antes de modelar y fingerprint
\texttt{c44de0783d3c3f88ef6151eaa75536d1020f5b208327419c4cb7f3ac459d3892}.

Por ausencia de GPU, los experimentos nuevos utilizaron backbones ImageNet
congelados y una cabeza logística entrenada sobre sus embeddings; no se
presentan como fine-tuning end-to-end. Optuna ejecutó 30 trials sobre
EfficientNet-B0 y podó 13. La mejor configuración elevó el macro-F1 de
validación de la cabeza baseline de 0.7799 a 0.8887, principalmente al usar
balanceo por raíz cuadrada, una tasa menor y penalización L2.
EfficientNet-B0 superó a ShuffleNetV2-x1.0, MobileNetV3 Small y FastViT-T8 como
modelo individual. Los errores fueron suficientemente distintos para que el
weighted soft voting mejorara hasta 0.9202 en validación, con pesos aprendidos
sin consultar el holdout.

En cinco folds, el ensemble obtuvo macro-F1 medio de 0.9097 y desviación 0.0090.
La evaluación única del holdout produjo 0.9537 de accuracy, 0.9012 de precision
macro, 0.9092 de recall macro, 0.9035 de macro-F1 y 0.9537 de F1 ponderado.
Potasio mejoró de forma orientativa desde el F1 histórico 0.62 hasta 0.6706,
pero siguió siendo la clase más débil: 28 de sus 36 errores terminaron en
nitrógeno. Fairness encontró las mayores brechas por clase (0.3206 de F1) y
fuente (0.3151); además, roya obtuvo recall 1.0 en 322 imágenes de laboratorio
y 0.3125 en solo 16 de campo, una señal importante aunque de soporte pequeño.

Grad-CAM y LIME se aplicaron a casos finales, incluidos un error de alta
confianza y una deficiencia nutricional. La demo ejecutable genera Top-3,
confianza, latencia y advertencia sobre una imagen real. También se verificó la
aplicación React Native con EfficientNet-Lite0 TFLite, inferencia offline y
rechazo fuera de dominio: typecheck y 276 pruebas aprobaron. El endpoint y el
APK no estuvieron disponibles en la comprobación final, por lo que el
prototipo es funcional e integrado, pero no recibe puntuación completa de
despliegue público.

\noindent\textbf{Palabras clave:} clasificación foliar, maíz, Optuna, ensemble,
validación cruzada, fairness, interpretabilidad, inferencia móvil.
